\documentclass[11pt,a4paper]{article}

% === PACKAGES ===
\usepackage[utf8]{inputenc}
\usepackage[margin=1in]{geometry}
\usepackage{amsmath,amssymb,amsthm}
\usepackage{enumitem}
\usepackage{graphicx}
\usepackage{hyperref}
\usepackage{xcolor}
\usepackage[most]{tcolorbox}
\usepackage{fancyhdr}
\usepackage{bm}
\usepackage{tikz}

% === THEOREM ENVIRONMENTS ===
\theoremstyle{definition}
\newtheorem{theorem}{Theorem}[section]
\newtheorem{lemma}[theorem]{Lemma}
\newtheorem{proposition}[theorem]{Proposition}
\newtheorem{corollary}[theorem]{Corollary}
\newtheorem{definition}[theorem]{Definition}
\newtheorem{example}[theorem]{Example}
\newtheorem{remark}[theorem]{Remark}

% === CUSTOM SIDEBARS — Subtle left-border style ===
\newtcolorbox{professorsnote}{
    blanker,
    borderline west={2.5pt}{0pt}{blue!50!black},
    left=10pt,
    before skip=10pt,
    after skip=10pt,
    before upper={\textcolor{blue!50!black}{\textbf{Professor's Note.}}\enspace}
}

\newtcolorbox{insight}{
    blanker,
    borderline west={2.5pt}{0pt}{green!50!black},
    left=10pt,
    before skip=10pt,
    after skip=10pt,
    before upper={\textcolor{green!50!black}{\textbf{Insight.}}\enspace}
}

\newtcolorbox{intuition}{
    blanker,
    borderline west={2.5pt}{0pt}{orange!60!black},
    left=10pt,
    before skip=10pt,
    after skip=10pt,
    before upper={\textcolor{orange!60!black}{\textbf{Intuition.}}\enspace}
}

\newtcolorbox{interview}{
    blanker,
    borderline west={2.5pt}{0pt}{red!50!black},
    left=10pt,
    before skip=10pt,
    after skip=10pt,
    before upper={\textcolor{red!50!black}{\textbf{Interview.}}\enspace}
}

% === DOCUMENT INFO (CUSTOMIZE THESE) ===
\newcommand{\lecturetitle}{[LECTURE TITLE]}
\newcommand{\lecturenum}{[LECTURE NUMBER]}
\newcommand{\coursenum}{[COURSE NUMBER]}
\newcommand{\coursename}{[COURSE NAME]}
\newcommand{\semester}{[SEMESTER]}
\newcommand{\lecturedate}{[DATE]}

% === HEADER/FOOTER ===
\pagestyle{fancy}
\fancyhf{}
\rhead{Lecture \lecturenum}
\lhead{\lecturetitle}
\rfoot{Page \thepage}

% === COMMON MATH OPERATORS (add more as needed) ===
\DeclareMathOperator{\E}{\mathbb{E}}
\DeclareMathOperator{\Var}{Var}
\DeclareMathOperator{\Cov}{Cov}
\DeclareMathOperator{\tr}{tr}
\DeclareMathOperator{\rank}{rank}
\DeclareMathOperator*{\argmax}{arg\,max}
\DeclareMathOperator*{\argmin}{arg\,min}
\newcommand{\R}{\mathbb{R}}
\newcommand{\N}{\mathbb{N}}
\newcommand{\Z}{\mathbb{Z}}
\newcommand{\C}{\mathbb{C}}

% === TITLE ===
\title{\textbf{\coursename} (\coursenum) \\ Lecture \lecturenum: \lecturetitle}
\author{Lecture Notes}
\date{\lecturedate \\ \semester}

\begin{document}

\maketitle

\tableofcontents
\newpage

% === LECTURE CONTENT GOES HERE ===
% 
% Structure by topic, not by slide number.
% Use theorem environments for formal statements.
% Use sidebars for commentary:
%   - \begin{professorsnote}...\end{professorsnote} for professor's opinions/advice
%   - \begin{insight}...\end{insight} for your added explanations
%   - \begin{intuition}...\end{intuition} for proof strategies and "why it works"
%   - \begin{interview}...\end{interview} for quant interview connections
%
% Every equation should be grammatically integrated into prose.
% Define notation before first use.
% Motivation precedes formal statements.

\section{Introduction}

[Overview of today's lecture topic and how it connects to previous material.]

\section{Main Content}

[Organized by logical topic flow, not slide order.]

\subsection{Topic 1}

[Content here...]

\begin{definition}
[Formal definition from slides.]
\end{definition}

\begin{professorsnote}
[What the professor said about this — their emphasis, practical advice, etc.]
\end{professorsnote}

\begin{theorem}
[Formal theorem statement from slides.]
\end{theorem}

\begin{proof}
[Proof as given in lecture — from audio/board work, not textbook version.]
\end{proof}

\begin{intuition}
[Why this theorem makes sense, proof strategy explained.]
\end{intuition}

\subsection{Topic 2}

[Continue with logical flow...]

\section{Summary}

[Key takeaways from today's lecture.]

\end{document}
